\documentclass{handout}

\assignment{Homework 7}
\title{Electric and magnetic fields}
\duedate{Wednesday November 5, 2025}

\begin{document}
\maketitle

\hwinstructions

\begin{questions}
    \question \textbf{Potential fields}: For the following potential field, mark where there are regions of positive and negative charge. Explain your reasoning in terms of gradients, the electric field, and divergence.
    \begin{figure}[h!]
        \centering
        \includegraphics[height=4in]{Images/Peaks.png}
    \end{figure}
    \begin{figure}[h!]
        \centering
        \includegraphics[height=2in]{Images/GaussDice.pdf}
    \end{figure}

    \question \textbf{Charge on a die}: You come across the following die, with its faces numbered 1-6 in the usual way (opposite faces add to 7). Upon measuring the electric field, you find that at all locations the electric field points to the right (in the $+y$ direction) with varying strengths. On the left of the die it has strength $|\vec{E}|=E_{3}$, on the right $|\vec{E}|=E_{4}$, and on all other faces $|\vec{E}|=E_{other}$. The die is 1.5 cm on each side.
    \begin{parts}
        \part Write an expression ($\pm \hat{x}, \pm \hat{y}$, etc.) for the outward pointing normal vector of each face (1-6).
        \part Write expressions for the dot product $(\vec{E} \cdot \hat{n})$ for each face (hint: some faces should be zero)
        \part Write expressions for the electric flux over each face.
        \part If $E_3=1.2\times 10^6$ N/C, $E_{other}=1.8\times 10^6$ N/C, and $E_{other}=2.5\times 10^6$ N/C, how much charge is inside the cube? (Remember $\Phi_E=Q_{in}/\epsilon_0$ where $\epsilon_0= 8.854\times 10^{-12}$ F/m.
    \end{parts}

    \question \textbf{Cross products}: Compute the cross products of the following vectors. For each, compute using component form and find the angle between the vectors. Also, draw a picture and explain how the right hand rule applies to your drawing.
    \begin{parts}
        \part $\langle -2,3, 0\rangle \times \langle 3,0,0\rangle$
        \part $\langle -2,-3, 0\rangle \times \langle 3,-4,0\rangle$
        \part $\langle 1,1, 0\rangle \times \langle 0,0,1\rangle$
        \part $\langle 2,1, 0\rangle \times \langle 0,0,-1\rangle$
    \end{parts}

    \question \textbf{Magnetic force}: A proton (mass $1.66\times 10^{24}$ kg, charge $+1.6\times 10^{-19}$ C), is moving with a velocity of $\langle-3,2,0\rangle$ km/s in the neighborhood of a magnetic field with value $\langle 2,4,0\rangle\times 10^{-4} $ T.
    \begin{parts}
        \part What is the magnitude of the force on the proton?
        \part Draw the vectors in question and explain their directions using the right-hand rule.
    \end{parts}

    \question \textbf{Cyclotron motion}: You observe a proton (see previous for properties) moving in a magnetic field of strength 3 G ($3\times 10^{-4}$ T) out of the page (towards your face). It is moving in a circular path with a 5 cm radius.
    \begin{parts}
        \part Is the proton moving clockwise or counterclockwise around this circle? What if it was an antiproton (same charge but negative, same mass).
        \part What is the speed of the proton?
    \end{parts}
\end{questions}

\end{document}
